% ECONOMICS_PROBLEM: {"id": "D91-637", "title": "Stability of reference dependence", "jel_code": "D91", "source_id": "OP-0062"}
% !TeX program = xelatex
\documentclass[11pt,a4paper]{article}
\usepackage[margin=23mm]{geometry}
\usepackage{fontspec}
\setmainfont{Latin Modern Roman}
\usepackage{amsmath,amssymb,mathtools}
\usepackage{xcolor,enumitem,booktabs,longtable,array,xurl,hyperref}
\definecolor{muted}{HTML}{53636C}
\hypersetup{colorlinks=true,urlcolor=blue,linkcolor=blue}
\setlength{\parindent}{0pt}
\setlength{\parskip}{5pt}
\newcommand{\R}{\mathbb R}
\newcommand{\E}{\mathbb E}
\newcommand{\N}{\mathbb N}
\newcommand{\ind}{\mathbf 1}
\newcommand{\Var}{\operatorname{Var}}
\newcommand{\Cov}{\operatorname{Cov}}
\newcommand{\supp}{\operatorname{supp}}
\DeclareMathOperator*{\argmin}{arg\,min}
\DeclareMathOperator*{\argmax}{arg\,max}
\newcommand{\entrymeta}[3]{{\sffamily\small\color{muted}\textbf{\texttt{#1}}\quad|\quad #2\par #3\par}}
\newcommand{\Problemref}[1]{\texttt{#1}}
\newcommand{\JELref}[2]{\texttt{#2} (JEL \texttt{#1})}
\begin{document}
\section*{Stability of reference dependence}
\textbf{Problem ID:} \texttt{D91-637}\quad \textbf{JEL:} \texttt{D91}\par
% BEGIN ECONOMICS STATEMENT
\label{problem:OP-0062}
\label{atlas:prob:B2}

\noindent\textbf{Original identifier:} B2 (atlas item 62). \textbf{Classification:} Behavioral identification and stability research program. \textbf{Original tractability label:} Medium (editorial, not a complexity result).

\subsection*{Definitions and assumptions}

For individual $i$, domain $d$ in a specified finite collection, and occasion $t$, let $x$ be a scalar consumption outcome expressed in a specified domain-specific utility scale. Let $r_{idt}$ denote a deterministic reference point and consider the illustrative utility
\[
U_{idt}(x)=u_d(x)+\eta_{id}\,v_{\lambda_{id}}(x-r_{idt}),\quad
v_\lambda(z)=\begin{cases}z,&z\geq0,\\ \lambda z,&z<0,\end{cases}
\]
where $\eta_{id}\geq0$ and $\lambda_{id}\geq1$. Specify $u_d$, choice noise, information sets, and scale normalization. Expectations-based reference distributions require integration of gain--loss utility over that distribution; they cannot generally be replaced by its mean.

\subsection*{Formal research question}

Given longitudinal choices and randomized changes to status quo, elicited expectations, or goals, identify the reference map $r_{idt}=g_d(H_{it},Z_{idt};\gamma_d)$ and the joint distribution of $(\eta_{id},\lambda_{id})$. Here $H_{it}$ is observed history and $Z_{idt}$ is randomized reference manipulation. Characterize the experimental variation and exclusions sufficient to distinguish changes in $g_d$ from changes in curvature or choice noise. For a declared equivalence tolerance $\varepsilon>0$, assess whether cross-domain normalized parameters satisfy
\[
\Pr\!\left(\max_{d,d'}|\lambda_{id}-\lambda_{id'}|\leq\varepsilon\right)\geq1-\alpha,
\]
where $\alpha\in(0,1)$ is a specified population exception rate.

\subsection*{Interpretation and scope}

Reference dependence is a family of models, not a uniquely identified primitive inferred from any observed endowment effect. Monetary and nonmonetary domains need comparable normalizations before coefficient equality has content. Randomly assigning a possession may change information, transaction costs, or beliefs about quality, violating an exclusion that it changes only the reference point. Designs should elicit expectations before choice, avoid treating endogenous expectations as instruments, and allow reference points to adapt over time. Stability of individual rankings is weaker than stability of cardinal parameters. An empirical contribution could produce a confidence set for the fraction of individuals satisfying the displayed tolerance condition, or demonstrate observational equivalence under existing protocols. The model shown here intentionally isolates reference dependence; it is not the complete prospect-theory treatment of probability weighting.

\subsection*{Source audit and status}

The three references are verified foundational models: [S1] addresses risky choice, [S2] riskless reference dependence, and [S3] expectations-based reference points. Their distinct constructions should not be collapsed into one formula. They establish useful representations and predictions, not a universal invariant loss-aversion coefficient. The formal stability and identification targets above are editorial extensions; the audit does not claim that these old publications certify their open status as of 2026.

\subsection*{References}

\begin{enumerate}[label={[S\arabic*]},leftmargin=*]

\item Daniel Kahneman and Amos Tversky (1979). \emph{Prospect Theory: An Analysis of Decision under Risk}. Econometrica 47(2), 263--291. \url{https://doi.org/10.2307/1914185}\par \textit{Source role:} Foundational descriptive model. \textit{Locator:} Author bibliography; original model of risky choice.

\item Amos Tversky and Daniel Kahneman (1991). \emph{Loss Aversion in Riskless Choice: A Reference-Dependent Model}. Quarterly Journal of Economics 106(4), 1039--1061. \url{https://doi.org/10.2307/2937956}\par \textit{Source role:} Model of reference-dependent riskless choice. \textit{Locator:} Abstract and article metadata.

\item Botond Kőszegi and Matthew Rabin (2006). \emph{A Model of Reference-Dependent Preferences}. Quarterly Journal of Economics 121(4), 1133--1165. \url{https://doi.org/10.1093/qje/121.4.1133}\par \textit{Source role:} Endogenous expectations-based reference point. \textit{Locator:} Abstract and article metadata.

\end{enumerate}

\subsection*{Common conventions: Source mathematical conventions}

Unless an entry states otherwise, agent and firm sets are finite. State and action spaces are measurable, strategies and outcome maps are measurable, probability laws are specified, and expectations exist for the targets under discussion. Infinite-horizon discount factors lie in $(0,1)$ and discounted objectives are finite. Norms, tolerances and complexity input sizes must be specified for computational comparisons. Constants or primitives that appear locally are inputs, not empirical facts asserted by this catalogue.

\textbf{Identification.} A statistical model $m\in\mathcal M$ implies an observable law $P_m$ and target $\theta(m)$. At law $P$, its identified set is
\[
\Theta(P;\mathcal M)=\{\theta(m):m\in\mathcal M,\ P_m=P\}.
\]
Point identification holds when this set is a singleton, or equivalently when $P_{m_1}=P_{m_2}$ implies $\theta(m_1)=\theta(m_2)$ for all compatible models. Sharp bounds mean bounds attained, or approached when the boundary is not attained, by compatible models. Writing this set is a definition, not a solution: the substantive task is to establish informative restrictions and derive or estimate the set. An unrestricted-confounding model may give only trivial bounds.

\textbf{Inference.} Identification, estimability and valid uncertainty quantification are separate questions. Fix $\alpha\in(0,1)$. A confidence procedure $C_n$ over a declared law class $\mathcal P$ has asymptotic uniform coverage at level $1-\alpha$ when
\[
\liminf_{n\to\infty}\inf_{P\in\mathcal P}\Pr_P\{\theta(P)\in C_n\}\geq1-\alpha.
\]
For set-identified targets the coverage event must be defined separately, for example coverage of the full identified set. If an entry uses identification-indexed models instead of uniquely indexed laws, the same coverage convention is applied over those models. Pointwise limits do not establish uniform validity, and asymptotic validity does not guarantee finite-sample precision. Sample sizes, dependence structures and weak-identification sequences are part of each application.

\textbf{Counterfactuals.} $Y_i(a)$ is a potential outcome under a specified intervention, including its timing, implementation and versions. It is not $Y_i$ conditional on voluntarily choosing $a$. A full intervention vector permits interference; shorthand scalar treatment notation does not silently impose no interference. Assignment, selection, attrition, anticipation and measurement must be specified. A claim about an instrument's local effect is not automatically a population policy effect.

\textbf{Economic equilibrium.} A counterfactual model includes best responses, constraints, feasibility, market clearing, state transitions and expectations consistent with the stated information. When several equilibria exist, either a declared selector is an input or the target is set-valued. Comparisons across policy regimes must use the stated selection convention. Reduced-form relationships are not presumed invariant after an equilibrium-changing intervention.

\textbf{Optimization and welfare.} Optimization is over the stated feasible set. If attainment is not established, $\sup$ or $\inf$ and $\varepsilon$-optimal choices replace an asserted maximizer or minimizer. For example, with value $V=\sup_{a\in\mathcal A}W(a)<\infty$, an $\varepsilon$-optimal set is $\{a\in\mathcal A:W(a)\geq V-\varepsilon\}$ for $\varepsilon>0$. Benefits, costs, risk and health or environmental outcomes can be aggregated only using declared units and conversion weights. Interpersonal utility weights, the treatment of fiscal transfers and foreign profits, discounting, and the behavioral welfare criterion are explicit inputs. A comparative-static derivative additionally requires existence and appropriate differentiability of the selected counterfactual.

\textbf{Sources.} Local labels [S1], [S2], and so forth restart in every question. Locators identify the relevant abstract, model, section or result at the level actually checked; a bibliographic or abstract-level verification is not represented as inspection of an entire proof. Working papers are distinguished from later publications and changing versions. Source summaries are paraphrased. All URL checks and editorial formulations refer to the audit date stated above.
% END ECONOMICS STATEMENT
\end{document}
